\subsection{Kesimpulan}

Berdasarkan hasil perancangan, implementasi, dan pengujian aplikasi Jatidiri berbasis mobile menggunakan framework Flutter dengan metode Prototype, maka dapat disimpulkan beberapa hal sebagai berikut:

1. Aplikasi Jatidiri berhasil dirancang dan diimplementasikan sebagai aplikasi mobile berbasis Flutter yang menyajikan berbagai tes psikologi, meliputi tes kecerdasan, kendali diri, bakat, kebahagiaan, kecemasan, gaya belajar, kendali stres, kesehatan psikologis, potensi, karier, kecerdasan dewasa, dan ketangguhan.

2. Implementasi desain antarmuka pengguna (UI) dari Figma ke dalam Flutter dapat dilakukan dengan baik, sehingga menghasilkan tampilan aplikasi yang interaktif, responsif, dan sesuai dengan rancangan awal, meliputi halaman Splash Screen, Walkthrough, Register, Login, Forgot Password, Survei, Home, Daftar Tes, Detail Tes, Halaman Tes, dan Komunitas.

3. Penerapan metode Prototype terbukti efektif dalam mendukung proses pengembangan aplikasi secara iteratif, terutama dalam menyempurnakan tampilan dan alur navigasi aplikasi berdasarkan evaluasi berkelanjutan terhadap desain dan implementasi UI.

4. Berdasarkan hasil pengujian Black Box, seluruh halaman dan fungsi navigasi utama aplikasi berjalan sesuai dengan rancangan yang telah ditetapkan. Karena aplikasi masih berupa prototype tanpa backend, pengujian difokuskan pada aspek tampilan antarmuka (UI), navigasi, dan interaksi pengguna, yang semuanya dinyatakan berfungsi dengan baik.


\subsection{Saran}

Berdasarkan keterbatasan dan hasil penelitian yang telah dilakukan, maka disarankan beberapa hal untuk pengembangan selanjutnya, yaitu:

1. Mengembangkan aplikasi dengan menambahkan backend (server dan database) agar fitur register, login, penyimpanan hasil tes, dan riwayat pengguna dapat berfungsi secara nyata.

2. Mengintegrasikan sistem autentikasi pengguna yang aman, misalnya menggunakan Firebase Authentication atau REST API berbasis backend.

3. Menambahkan fitur perhitungan dan analisis hasil tes psikologi secara otomatis serta menampilkan interpretasi hasil kepada pengguna.

4. Melakukan pengujian usability dengan melibatkan pengguna nyata untuk mengevaluasi pengalaman pengguna (UX) secara lebih mendalam.

5. Mengembangkan aplikasi agar tidak hanya berjalan di Android, tetapi juga dapat diimplementasikan pada iOS dan web.

6. Menyempurnakan tampilan UI/UX berdasarkan umpan balik pengguna dan standar desain terkini.
